\documentclass[10pt,a4paper]{amsart}
\usepackage[margin=19mm]{geometry}
\usepackage{amsmath,amssymb,mathtools}
\usepackage[scr=rsfs,cal=euler]{mathalfa}
\usepackage{xfrac}
\usepackage[unicode,hidelinks,bookmarks=false]{hyperref}
\newcommand{\ie}{i.e.\ }
\newcommand{\cf}{cf.\ }
\newcommand{\resp}{respectively\ }
\newcommand{\Subsection}{\subsection}
\providecommand{\nobreakdash}{\nobreak\hbox{-}}
\setlength{\emergencystretch}{2em}
\multlinegap0pt
\usepackage{bm}
\usepackage{bbm}
\usepackage{dsfont}
\usepackage{xspace}
\usepackage{cancel}
\usepackage{mathtools}
\usepackage{amsthm}
\makeatletter
\@ifundefined{theorem}{\newtheorem{theorem}{Theorem}}{}
\@ifundefined{lemma}{\newtheorem{lemma}[theorem]{Lemma}}{}
\@ifundefined{proposition}{\newtheorem{proposition}[theorem]{Proposition}}{}
\makeatother
\newcommand{\calM}{\mathcal{M}}
\title[Partial fraction decompositions on hyperplane arrangements]{Partial fraction decompositions on hyperplane arrangements}
\author{Claire de Korte, Teresa Yu}
\date{Source: arXiv:2602.06531v2 (CC BY 4.0)}
\begin{document}
\maketitle

\noindent\emph{Source and licence.} Partial fraction decompositions on hyperplane arrangements. Version arXiv:2602.06531v2 (CC BY 4.0), distributed under the Creative Commons Attribution 4.0 International licence, as recorded in the arXiv licence metadata for this exact version and shown on the abstract page https://arxiv.org/abs/2602.06531v2. Full rights notice: licenses/arxiv-2602.06531-CC-BY-4.0.txt.\par\smallskip

\noindent\textit{Editorial scope.} Counted unit: the Proposition whose source label is \texttt{prop: linear independence of products} in the article (source0.tex lines 216--222, source label \texttt{prop: linear independence of products}), delivered together with the article's own complete original proof of it (source0.tex lines 224--244, one \texttt{proof} environment of 21 lines). No mathematics is written, repaired or re-derived: statement and proof are the article's own text line for line. Required context retained uncounted: lem:LT-generic (lines 246--251). Every definition, display and label of those lines is retained unchanged. Omitted from this package and not redistributed: 1--215, 223--223, 245--245, 252--1040. Nothing inside a retained excerpt is omitted or altered: every retained body is delivered exactly as the article's lines are written, so no omission receipt of any kind accompanies this package. Citations: the retained excerpts carry no citation command at all, so no bibliography entry of the article is transcribed and no reference is redistributed. Reference adaptation: none was needed and none was made; every reference inside the delivered unit resolves inside it, so no unresolved reference or citation is produced. Printed slips kept literal: no printed word, symbol, hypothesis or number of the counted statement or of its proof was altered. Layout and class: the article's own class is \texttt{amsart} with packages including babel, fullpage, amsmath, amsfonts, amsthm, graphicx, hyperref, mathtools, comment, cleveref, bbm, nicematrix, tikz, array; the wrapper is the lane's pinned \texttt{amsart} editorial wrapper at 10pt on A4, which restates the theorem-like environments the retained text needs and adds the article's own macro definitions that the retained text needs (\texttt{\textbackslash calM}), with extra wrapper packages (bm, bbm, dsfont, xspace, cancel, mathtools). The delivered numbering is the unit's own. Assets: no retained span contains a figure environment, a graphics-inclusion command, a PDF-page inclusion, a table or a TikZ picture; no image file is copied into this package, the package declares no asset dependency and no image file is redistributed with it. Licence: version arXiv:2602.06531v2 of this article is distributed under the Creative Commons Attribution 4.0 International licence, as recorded in the arXiv licence metadata for this exact version and shown on the abstract page https://arxiv.org/abs/2602.06531v2. A full rights notice is delivered at licenses/arxiv-2602.06531-CC-BY-4.0.txt. Characters: every retained excerpt is pure ASCII, so the delivered engine, XeLaTeX with its default Latin Modern text font, reports no missing character.
\begin{lemma}\label{lem:LT-generic}
    Following the notation of the proof of Proposition \ref{prop: linear independence of products}, we have that $h=\pm 1$:
    \begin{equation}\label{eqn:lem-LT}
        \det(\mathcal{M})=\pm\prod_{1\leq i_i<\cdots<i_r\leq n}A[i_1,\ldots, i_r].
    \end{equation}\
\end{lemma}

\begin{proposition}\label{prop: linear independence of products}
With the notation set above, we have:
\begin{equation}
\det(\mathcal{M})=\pm\prod_{1\leq i_1<\cdots<i_r\leq n}A[i_1,\ldots, i_r]. 
\end{equation}
In particular, if every $r$-subset of the linear forms $\ell_1,\ldots,\ell_n$ is linearly independent, then the set of $d$-fold products of the linear forms is also linearly independent.
\end{proposition}

\begin{proof}
 We may assume that $K$ is algebraically closed. Note that each \(A[i_1,\ldots, i_r]\) is irreducible, and so by Hilbert's Nullstellensatz,
\[A[i_1,\ldots, i_r]\text{ divides } \det(\mathcal{M}) \ \Leftrightarrow \ \mathcal{V}(A[i_1,\ldots, i_r])\subset \mathcal{V}(\det(\mathcal{M})).\]

We now show that $\mathcal{V}(A[i_1,\ldots, i_r])\subset \mathcal{V}(\det(\mathcal{M}))$. Suppose that there exists some relation between the \(\ell_{i}\), i.e., for some \(\{i_1,\ldots,i_r\}\subset[n]\), there exists \(c_1,\ldots, c_r\in K\) not all zero such that:
\begin{equation}\label{eq: linear dependence relation}
\sum_{j=1}^{r}c_i\ell_{i_j}=0.
\end{equation}
The existence of such a relation is equivalent to \(A[i_1,\ldots, i_r]\) being identically zero. In other words, the tuples \((a_{11},\ldots,a_{nr})\in K^{n\times r}\) for which such a relation (\ref{eq: linear dependence relation}) exist form \(\mathcal{V}(A[i_1,\ldots, i_r])\).
Define the set \(\{i'_1,\ldots,i'_{d-1}\}=[n]\setminus\{i_1,\ldots,i_r\}\). Multiplying each term in (\ref{eq: linear dependence relation}) by \(\ell_{i'_1}\cdots\ell_{i'_{d-1}}\), we find that:
\begin{equation}\label{eq: product linear independence relation}
\sum_{j=1}^{r}c_i\ell_{i_j}\ell_{i'_1}\cdots\ell_{i'_{d-1}}=0.
\end{equation}
Each \(c_i\) in (\ref{eq: product linear independence relation}) is multiplied by a unique collection of \(d\) linear forms.
But then we see that the set of $d$-fold products of the $\ell_i$ is linearly dependent, and so $\det(\calM)$ vanishes on \(\mathcal{V}(A[i_1,\ldots, i_r])\). Thus \(\mathcal{V}(A[i_1,\ldots, i_r])\subset \mathcal{V}(\det(\mathcal{M}))\). This implies that:
\[\det(\mathcal{M})=h\cdot\prod_{\substack{\{i_1,\ldots, i_r\}\subset[n] \\ i_i<\cdots<i_r}}A[i_1,\ldots, i_r],\qquad h\in K[a_{11},\ldots,a_{nr}].\]

Each $A[i_1,\ldots, i_r]$ is a degree $r$ polynomial, and so their product has degree $\binom{n}{r}\cdot r$ in the $a_{ij}$.
The entries of the matrix \(\mathcal{M}\) are each of degree \(d\) in the \(a_{ij}\), and so \(\det(\mathcal{M})\) has degree \(\binom{n}{d}\cdot d\). For \(n=r+d-1\), we have \(\binom{n}{r}\cdot r=\binom{n}{d}\cdot d\), so \(\det(\mathcal{M})\) and the product of $A[i_1,\ldots, i_r]$ are polynomials of the same degree. This is only possible if \(h\in K\). The fact that $h=\pm 1$ is proven in the
following lemma (Lemma \ref{lem:LT-generic}).
\end{proof}

\end{document}
